% ===== tr-05 块2b：第 4 章 pp.44--47（原文页 44--47）=====
% [原文 p.44]
这里，幺半群 $Q$ 由「该曲线是基本的」这一事实所决定。我们对此作热带化，于是得到图表
\begin{center}
\begin{tikzcd}
\operatorname{Trop}(C^\dagger) \arrow[r,"\operatorname{Trop}(f)"] \arrow[d,"\operatorname{Trop}(\pi)"] & \operatorname{Trop}(\mathcal{X}_0^\dagger) \arrow[d,"\operatorname{Trop}(q)"] \\
\operatorname{Trop}(W^\dagger)=\sigma_Q^\vee \arrow[r,"\operatorname{Trop}(g)"] & \operatorname{Trop}(0^\dagger)=\mathbb{R}_{\ge 0}
\end{tikzcd}
\end{center}

$\operatorname{Trop}(\pi)$ 的纤维一般是一维图，而 $\operatorname{Trop}(q)^{-1}(1)$ 是 $\mathcal{X}_0^\dagger$ 的对偶相交复形 $B$。（一般地，它只是一个多面体复形，不像环面退化的情形那样在余维一上带有仿射结构。）于是 $\operatorname{Trop}(g)^{-1}(1)\subseteq\sigma_Q^\vee$ 可以看作参数化「从图（$\operatorname{Trop}(\pi)$ 的纤维）到 $B$ 的映射」的空间。这些将是热带曲线。事实上，在 $B$ 确实带有仿射结构的地方，这些曲线满足热带平衡条件。

与基本对数结构相伴的幺半群 $Q$ 必须满足的根本性质是：$\operatorname{Trop}(g)^{-1}(1)$ 必须参数化 $B$ 中所有具有同一「组合类型」的热带曲线。这就精确刻画了热带曲线与对数曲线之间的对应。

我们还可以描述粘合公式的预期形态，它与 Brett Parker 在其设定中发展的公式 [67] 相呼应。如上考虑 $B$ 中的热带曲线。这些曲线一般成族移动，但对任意给定的数据 $\beta$，总会有有限多条代表 $\beta$ 的热带曲线，它们在不改变定义域图的情况下无法形变。我们称这样的热带曲线为\emph{刚性的}。于是实际的模空间 $\mathcal{M}(\mathcal{X}_0^\dagger/0^\dagger,\beta)$ 可以看作具有一个由这些刚性曲线标号的「虚拟不可约分支分解」。此外，与任一刚性曲线相伴的「虚拟不可约分支」还可以进一步与热带曲线的每个顶点所对应的曲线模空间联系起来。这最终应当使我们能够把 $\mathcal{X}_0^\dagger/0^\dagger$ 的 Gromov--Witten 不变量、从而把 $\mathcal{X}_0^\dagger$ 的某个光滑化的 Gromov--Witten 不变量，用简单得多的不变量表达出来。这是与 Abramovich、Chen 和 Siebert 正在合作进行中的项目。

\section{B-模型与热带几何}
让我们转向 $B$-模型，并理解热带几何如何可能在 $B$-模型计算所需的复结构变分中显现出来。

这个问题与重构问题（如问题 4.7.12 所述）密切相关。若给定 $(B,\mathcal{P},\varphi)$，人们能够找到一个环面退化 $\mathcal{X}\to D$ 的显式描述，使其对偶相交复形为 $(B,\mathcal{P},\varphi)$，那么就可以利用这一显式描述计算周期，并提取出镜像的 $B$-模型预言。
% [原文 p.45]
% [Tutor Note: 问题 4.5.6 指原文 §5 中陈述的重构问题第一版；p.35 的问题 4.7.12 原题标为「The reconstruction problem, Version II」，两处是同一问题的不同版本，不可混为同一编号。]
在描述这个问题的解决方案之前，让我先讲一点重构问题的历史。如问题 4.5.6 所述的那个版本，最早由 Fukaya 在 [13] 中研究。他在那里直接考虑扰动 $X_\epsilon(B_0)$ 上的复结构的问题，方法是考察支配复结构形变的 Kodaira--Spencer 方程。在 $\dim B=2$ 的情形下作非正式的论证时，他提出扰动应当集中在由梯度流线构成的树上，而这些流线最初从 $B$ 的奇点发出。这给出了第一个提示：重构问题的一个好解法或许真的会看到与曲线有关的东西。然而，Fukaya 的工作没有给出确定的定理，而且其分析看起来很可能非常困难。

2004 年，Siebert 与我在考虑如何用我们的纲领求解重构问题。给定 $(B,\mathcal{P})$，我们已在 [31] 中说明如何构造对数概形 $X_0(B,\mathcal{P},s)^\dagger$，连同到 $0^\dagger$ 的态射，它具有人们对环面退化的中心纤维 $\mathcal{X}_0^\dagger\to 0^\dagger$ 所期望的一切性质。我们最初的希望是：Bogomolov--Tian--Todorov 无障碍性定理的某种推广能让我们证明这类对数概形可以光滑化。特别地，Kawamata 与 Namikawa [51] 在正规交叉情形下用这一观点取得了成功。虽然这一进路在维数 $2$ 时很容易奏效，我们却无法让它在高维奏效。此外，这一进路未能给出描述 $B$-模型所需的、关于光滑化的显式描述。于是我们转向一种更显式的进路，即把显式的局部模型粘合起来。

当我们正在研究这一进路时，Kontsevich 与 Soibelman 在 [54] 中用刚性解析流形取代复流形，绕过了 Fukaya 进路中困难的分析。他们能够证明：给定一个带有 $24$ 个 focus-focus 型奇点（这是仿射曲面中出现的最简单的奇点类型，稍后将予描述）的热带仿射曲面 $B$，就可以构造一个刚性解析 K3 曲面 $X^{an}(B)$。这是通过用附着在 $B$ 上一些直线上的自同构把标准片段粘合起来而做到的。这些直线由梯度流线给出，从而得到一个与 Fukaya 所给出的相似、但精确得多的图景。

把我们粘合局部模型的进路与 Kontsevich--Soibelman 工作 [54] 的核心想法之一结合起来，我们便能在所有维数上完成这一构造，从而在代数几何范围内给出重构问题的一个令人满意的解。这项工作在 [33] 中完成。

% [Tutor Note: 原文用带 check 的 $\check X(B)$ 表示镜像侧的流形、用 $X(B)$ 表示非镜像侧；紧接着的一段中二者交替出现（$\check X(B)$ 上放辛结构、$X(B)$ 上讨论复结构），抄写时勿把 check 漏掉或加错对象。]
在综述这一进路之前，让我作一个哲学性的评论。注意，当我们讨论 A-模型时，我们曾观察到 $B$ 上的热带曲线应当对应于 $X(B)$ 上的全纯曲线。如果我们想让这些同样的热带曲线在镜像的 $B$-模型上也起作用，那么我们就不应当在形如 $X(B)$ 的复流形上、而应当在 $\check X(B)$ 上考虑 $B$-模型。这是一种略令人困惑的角色互换。通常计数曲线是在辛范畴中进行的，此处它应当意味着在
% [原文 p.46]
辛流形 $\check X(B)$ 上进行，而任何与复结构有关的事情都应当在 $X(B)$ 上做。这一角色互换可以如下解释。若我们要研究 $\check X(B)$ 上的伪全纯曲线，就需要在 $\check X(B)$ 上放一个殆复结构。做到这一点的一种办法是在 $B$ 上选取一个度量；这诱导出 $\check X(B)$ 上一个在环面纤维化各纤维上为常值的殆复结构，推广了第 2 节中由 Hessian 度量构造复结构的做法。于是在一个把殆复结构重新缩放的适当绝热极限中，伪全纯曲线预期会趋向于梯度流线构成的树。若所选度量事实上是 Hessian 的，则这些梯度流线事实上将是相对于 Legendre 对偶仿射结构的直线，因此这些梯度流线树可以看作热带曲线的推广。不过，热带几何是线性的，控制起来容易得多。我们采取的态度是：应当在热带几何出现的那一侧工作。事实上，这被证明是非常有帮助的。

有了这一点，我们就可以给出镜像对称纲领的一个略作修改的版本：
\begin{enumerate}
\item 从带有丰沛极化的 Calabi--Yau 流形环面退化 $\mathcal{X}\to D$ 出发。
\item 由这些数据构造对偶相交复形 $(B,\mathcal{P},\varphi)$。
\item 构造一个新的环面退化 $\check{\mathcal{X}}\to D$，使其\emph{相交复形}为 $(B,\mathcal{P},\varphi)$。这一退化应当由热带数据控制。
\item 用 $B$ 的热带几何来理解 $\mathcal{X}\to D$ 一般纤维上的亏格 $0$ 全纯曲线（或人们所关心的 A-模型的其他方面）。
\item 用 $B$ 的热带几何来理解 $\check{\mathcal{X}}\to D$ 的 Hodge 结构变分。
\item 利用 $\mathcal{X}$ 与 $\check{\mathcal{X}}$ 的 A-模型与 $B$-模型分别受 $B$ 上同一热带几何控制这一事实，来证明镜像对称。
\end{enumerate}

这里我们概述 [33] 中对步骤 (3) 的完成。第一步如下。给定 $(B,\mathcal{P},\varphi)$，我们希望构造退化的中心纤维 $\mathcal{X}_0^\dagger$。这其实已在 [31] 的 §5 中完成，其中对 $B$ 的奇异轨迹作了若干一般位置假设。作为概形，$\mathcal{X}_0$ 应当是什么是相当显然的。对每个极大胞腔 $\sigma$，有一个相伴的、以 $\sigma$ 为 Newton 多胞形的射影环面簇 $\mathbb{P}_\sigma$。任何面 $\tau\subseteq\sigma$ 指定一个环面分层 $\mathbb{P}_\tau\subseteq\mathbb{P}_\sigma$；并且给定 $\sigma_1,\sigma_2$ 极大时的 $\tau=\sigma_1\cap\sigma_2$，我们可以按环面等变的方式把环面分层 $\mathbb{P}_\tau\subseteq\mathbb{P}_{\sigma_1},\mathbb{P}_{\sigma_2}$ 粘合起来。当然，可能的粘合有一整族，由我们在 [31] 中称为\emph{闭粘合数据}的东西参数化。给定闭粘合数据 $s$，我们得到一个概形 $\check X_0(B,\mathcal{P},s)$。

现在，$\check X_0(B,\mathcal{P},s)$ 不可能成为环面退化的中心纤维，除非它带有正确类型的对数结构。这有许多可能不发生的原因。若 $s$
% [原文 p.47]
% [Tutor Note: 原文写 $\coprod_{v\in\mathcal{P}}V(v)$，这里的下标 $v\in\mathcal{P}$ 实际指对 $\mathcal{P}$ 的\emph{顶点} $v$ 取不交并（$\mathcal{P}$ 同时表示多面体分解及其顶点集），并非对 $\mathcal{P}$ 的所有元素取并。]
取得不好，则 $\check X_0(B,\mathcal{P},s)$ 可能存在零维分层，其邻域并不局部 étale 同构于某个仿射环面簇的环面边界；这是一个最低限度的前提。因此，我们必须把注意力限制在由我们所谓的\emph{开粘合数据}诱导出的闭粘合数据上。具体地说，$\mathcal{P}$ 的每个顶点 $v$ 如下定义局部模型 $V(v)\subseteq U(v)$。分段线性函数 $\varphi$ 在局部上只定义到相差仿射线性函数为止。选取 $\varphi$ 在 $v$ 的邻域内的一个代表 $\varphi_v$，它在 $v$ 处取值 $0$。通过在每个胞腔上线性地延拓该函数，我们可以把它看作扇 $\Sigma_v$ 上的分段线性函数（把 $\Sigma_v$ 看作某个 $\mathbb{R}^n$ 中的扇）。于是我们可以令
\[
P_v:=\{(m,r)\in\mathbb{Z}^n\times\mathbb{Z}\mid r\ge\varphi_v(m)\}.
\]
注意到 $(0,1)\in P_v$，我们令
\[
U(v):=\operatorname{Spec}\mathbb{C}[P_v],\qquad
V(v):=\operatorname{Spec}\mathbb{C}[P_v]/(z^{(0,1)}).
\]
注意 $z^{(0,1)}$ 在 $U(v)$ 的每个环面除子上恰好一阶消失，故事实上 $V(v)$ 是 $U(v)$ 的环面边界。正如我们在 [31] 中所证明的，$\check X_0(B,\mathcal{P},s)$ 成为环面退化中心纤维的一个必要条件是：它由把 $\coprod_{v\in\mathcal{P}}V(v)$ 对一个等价关系作商而得到。换言之，我们是在沿 Zariski 开子集把诸 $V(v)$ 粘合起来，以得到一个概形。\footnote{[31] 允许 $\mathcal{P}$ 的胞腔自交的情形。因此该等价关系仅仅是 étale 的，得到的是一个代数空间。}同样，粘合有一定选择余地，但现在粘合数据由诸 $V(v)$ 的开子集之间的等变等同给出。我们称之为开粘合数据。

采用开粘合数据的好处是：每个 $V(v)$ 都带有由除子型对数结构 $V(v)\subseteq U(v)$ 诱导的对数结构。这些对数结构在开粘合映射下并不被等同，但它们的幽灵层是同构的。于是幽灵层 $\overline{\mathcal{M}}_{V(v)}$ 粘合起来，给出一个幺半群幽灵层 $\overline{\mathcal{M}}_{\check X_0(B,\mathcal{P},s)}$。由此我们看到 $\varphi$ 如何影响对数结构。

接下来人们试图构造一个以该幽灵层为幽灵层的对数结构。这在 [31] 中是通过构造幽灵层与 $\mathcal{O}_{\check X_0(B,\mathcal{P},s)}^\times$ 的适当扩张来完成的，而这种扩张依赖于某些模（一般而言可能是空的）。好的情形是：可以找到复余维数至少为二的闭子集 $Z\subseteq\check X(B,\mathcal{P},s)$，以及 $\check X_0(B,\mathcal{P},s)$ 上的一个对数结构，连同态射 $\check X_0(B,\mathcal{P},s)^\dagger\to 0^\dagger$，后者在 $Z$ 之外对数光滑。此外，$\check X_0(B,\mathcal{P},s)\setminus Z$ 上的幽灵层应当是给定的幺半群幽灵层 $\overline{\mathcal{M}}_{\check X_0(B,\mathcal{P},s)}$ 限制到 $\check X_0(B,\mathcal{P},s)\setminus Z$ 上所得。我们称这样的带 $0^\dagger$ 上态射的对数概形为一个\emph{对数 Calabi--Yau 空间}。

[31] 的技术核心是对具有给定相交复形 $(B,\mathcal{P},\varphi)$ 的对数 Calabi--Yau 空间的显式分类，这是在模去关于 $B$ 的奇点的某些被称为单纯性（simplicity）的假设之下进行的。
